\section{Discusión}

\subsection{El modelo propuesto conserva el núcleo de Eriksen}

La diferencia estructural principal no consiste en sustituir la
respiración larvaria. Los términos $r_{C,m}$, $r_{C,B}$ y $r_{C,L}$
permanecen activos y continúan dependiendo de $m$, $Y_B$ y $Y_L$. El
modelo propuesto añade un estado microbiano y una dependencia de la
asimilación respecto del alimento disponible.

Por esta razón, los cambios de resultados entre estudios no pueden
explicarse únicamente por ``agregar el lecho''. En los datos propios
también cambian $a_{\max}$, $\rho$ y el factor dietario. Además, el
acople modifica la trayectoria larvaria a través de $D_C$, de modo que
el componente microbiano y el componente larvario dejan de ser
independientes.

\subsection{Interpretación del ajuste y de la validación interna}

En la calibración completa, el RMSE censurado conjunto disminuyó de
$183{,}61$ a $85{,}16$ ppm/min, equivalente a $53{,}6$ %. Este
resultado evidencia una mayor capacidad de ajuste del modelo completo
recalibrado, pero no constituye por sí solo evidencia predictiva fuera
de muestra.

En la validación cruzada interna agrupada por dieta y día, el modelo
acoplado redujo el RMSE censurado de $183{,}61$ a $124{,}94$ ppm/min,
una reducción puntual de $31{,}95$ % frente a Eriksen original. Sin
embargo, no mejoró el RMSE de los cierres exactos, que fue $137{,}20$
frente a $128{,}81$ ppm/min, ni la cobertura conjunta de las cotas
saturadas, que permaneció en $25$ %. Además, subestimó los dos cierres
de D1 en el día 6 y sobreestimó los cierres tardíos de D1 y D4.

El remuestreo bootstrap exploratorio confirmó una incertidumbre amplia:
los percentiles $2{,}5$ y $97{,}5$ de la reducción del RMSE censurado
frente a Eriksen fueron $-82{,}19$ y $52{,}45$ %, respectivamente.
Como este rango incluye reducciones negativas, la mejoría puntual no
debe interpretarse como superioridad predictiva concluyente. Tampoco
permite atribuir por separado los residuos restantes a una fuente
larvaria o microbiana específica.

\subsection{Ventajas y límites frente a estudios externos}

Los resultados externos impiden afirmar superioridad universal. En
Laganaro, la parametrización original de $Y_B$ y $m$ captura un efecto
de calidad dietaria que el factor de disponibilidad no reproduce. En
Bekker, la reducción del error individual convive con picos anticipados
y una partición microbiana casi invariante frente a la humedad.

En Ermolaev, el modelo propuesto mejora las tasas puntuales pero empeora
el sesgo del acumulado. En Schøn, el aumento monótono inducido por
$Q_{10}$ resulta incompatible con el colapso a alta temperatura. En
Rossi, una menor distancia agregada no equivale a reproducir la
variabilidad entre corridas ni el momento del pico.

Estas discrepancias muestran que un RMSE menor puede coexistir con una
dinámica o una partición incorrectas. La evaluación debe combinar
errores puntuales, acumulados, momentos de los picos y fracciones de
las fuentes.

\subsection{Descomposición controlada de la mejora}

La comparación entre Eriksen original, el núcleo recalibrado, el lecho
aditivo y el modelo acoplado permitió separar contablemente los
componentes de la reducción del error en calibración. La recalibración
del núcleo aportó $55{,}8$ % de la reducción total, la adición de
\co{} microbiano $12{,}7$ % y el acople con la disponibilidad de
alimento $31{,}5$ %. Por construcción, el lecho aditivo conserva la
trayectoria larvaria; en cambio, el acople la modifica mediante
$f_d\Phi_D(D_C)$.

Esta descomposición no identifica efectos causales independientes. La
recalibración modifica parámetros compartidos y el acople altera
simultáneamente la asimilación y las trayectorias de biomasa. Por tanto,
los porcentajes describen la secuencia de configuraciones evaluada y no
una partición única de las fuentes de error.

\subsection{Sensibilidad e identificabilidad}

Las perturbaciones locales de $\pm 10$ % mostraron que la suma de los
acumulados independientes de D1 y D4 fue más sensible a $a_{\max}$
($5{,}49$ %), $k_{\mathrm{ref}}(D1)$ ($3{,}06$ %) y
$\delta_{CS}(D4)$ ($2{,}75$ %). En contraste, las variaciones asociadas
con $\rho$ y $k_{\mathrm{ref}}(D4)$ fueron $0{,}18$ y $0{,}10$ %,
respectivamente. El parámetro $a_{\max}$ también produjo el mayor
cambio local del RMSE censurado, de $13{,}72$ ppm/min.

El perfil de $\rho$, reoptimizando $a_{\max}$ y $f_{D4}$, presentó una
región con aumento de costo no superior a $1$ % entre $4{,}3953$ y
$15$. Dentro de esa región, la suma de los acumulados permaneció entre
$237{,}06$ y $243{,}00$ g. Así, el límite superior de $\rho$ quedó mal
identificado, aunque el acumulado fue relativamente estable. El perfil
de $a_{\max}$, por su parte, favoreció valores próximos a su frontera
superior.

Estas regiones de costo casi equivalente son exploratorias: no son
intervalos de confianza ni intervalos predictivos. En consecuencia, la
estabilidad local del \co{} acumulado no equivale a una identificación
precisa de todos los parámetros ni elimina la incertidumbre predictiva
observada en la validación cruzada.
